\slajdid{09_1-s019}
\begin{frame}[label=09_1-s019]
\gusttekst\setlength{\abovedisplayskip}{5pt}\setlength{\belowdisplayskip}{5pt}
Zbog čega smo rekli da ostavljamo zavisnosti ulaznih i izlaznih napona u ovom obliku. Zato što nam je sada „lako“ da diferencirano i levu i desnu stranu po $V_I$
\[\frac{V_{DD}-V_O}{R_L}\approx\frac{k_n}{2}(V_I-V_{Tn})^2\quad\Bigg|\frac\partial{\partial V_I}\]
\[-\frac1{R_L}\frac{\partial V_O}{\partial V_I}\approx2\frac{k_n}{2}(V_I-V_{Tn})\]
pri čemu je
\[\frac{\partial V_O}{\partial V_I}=-1\]
\[\frac1{R_L}\approx k_n(V_I-V_{Tn})\]
\[V_{IL}\approx V_{Tn}+\frac1{k_nR_L}\]
\[V_{O(IL)}\approx V_{DD}-\frac1{2k_nR_L}\]
\end{frame}
